\subsection{Three existential normal forms after complex conjugation}

\begin{statement}{D23.C31}{Proved}{Complex conjugation identifies two cube-normalized cases}
\end{statement}
If a normalized actual mixed full-GT element exists, then
one exists whose degree-1344 data $(V,R)$ has one of the
three labels
\[
                           (1,0),\qquad(19,0),\qquad(27,2).
                                                               \tag{D23.C31.1}
\]
This uses cubing as in C30 and, if necessary, conjugation
by the actual element $c=(-1,1)$ of complex conjugation.
It is an existential normal form, not an assertion that
any of the three finite pairs lifts to an actual element.

\begin{proof}
First realize the complex-conjugation symmetry on the
chosen finite Nielsen orbit.  In
$G=\mathbb F_3^2\rtimes SL_2(\mathbb F_3)$, let
$d=\operatorname{diag}(1,-1)$.  The map
\[
            \eta(v,A)=(dv,dAd^{-1})
\]
is an involutive automorphism of $G$.  The exact
36-element order-six class used in P17 satisfies
$\eta(\mathcal C^{-1})=\mathcal C$, verified on every
class element in the certificate.

The standard GT element $c=(-1,1)$ sends each Artin
generator to its inverse.  The arithmetic identification
of this element is Lochak--Schneps,
\emph{Inventiones Mathematicae} \textbf{127} (1997),
571--600, Theorem 1 on printed page 573; the braid action
is the convention used in supplied FP.34 and D23.P18.
In the free-kernel
construction of P18 it therefore takes the four local
conjugacy classes to their inverse classes.  Postcompose
the target representation by $\eta$ to return to the
original class family.  The resulting bijection of the
finite Nielsen set intertwines the three braid
generators with their inverses.  The unique orbit of
degree 1344 in P17 must be preserved.  Hence it gives
a permutation $W$ with
\[
           WAW^{-1}=A^{-1},\quad WBW^{-1}=B^{-1},
           \quad WCW^{-1}=C^{-1}.                         \tag{D23.C31.2}
\]
The transitive labeled-graph calculation exhausts all
1344 possible images of one point and finds exactly
one such $W$; it also verifies $W^2=1$.

Postcomposition by a target-group automorphism commutes
with every precomposition defining a GT action.
Consequently the inserted $\eta$ cancels when one
conjugates a multiplier-one GT action by this symmetry:
the action of $chc^{-1}$ is $WV_hW^{-1}$.
On the normalized free-group action,
\[
              f_{chc^{-1}}(x,y)=f_h(x^{-1},y^{-1}).
\]
Since (D23.C31.2) inverts both $x=C^2$ and
$y=(ABA)^2$, its mixed value is $WR_hW^{-1}$.

The normalized characters remain in the required state.
On the canonical degree-36 detector the exact induced
$W$ sends the six first-pair labels by
\[
                       (\rho,a_3)\longmapsto(-\rho,a_3).
\]
This is checked on all six values.  For the marked
$S_4$ pair $p=(123)$, $t=(1234)$, conjugation by $(13)$
inverts both $p$ and $t$.  Thus the supplied value
$t^{2a_2}p^\rho$ goes to $t^{2a_2}p^{-\rho}$, so $a_2$
is unchanged as well.  In particular
$(\rho,a_3,a_2)=(0,1,0)$ is preserved.

Finally simultaneous conjugation of all eighteen
recorded pairs by $W$ preserves that list.  On the four
cube-normalized pairs of C30 it fixes $(1,0)$ and
$(27,2)$ and interchanges $(19,0)$ with $(23,1)$.
Conjugating the latter case if necessary gives exactly
the three normal forms (D23.C31.1).
\end{proof}

\paragraph{Certificate and scope.}
\path{agents/characters/code/hurwitz_complex_conjugation.py}
checks the inverse-class correction, reconstructs the
unique inverse-generator intertwiner, computes its
degree-36 detector action, verifies the marked $S_4$
inversion, and records its action on every candidate.
The JSON/NPZ records are
\path{agents/characters/records/hurwitz_complex_conjugation_probe.json}
and the NPZ of the same stem.  The actual-source inputs
are the standard $c=(-1,1)$ action, P18 functoriality,
and P17's unique degree-1344 orbit.  No arithmetic
lifting or extension of a finite candidate is assumed.
